% !PRINTED 182
\weberpage{182}
%<<CONT
it is no longer assumed that the determinant of the function $\varphi$
vanishes. This discriminant is
\weq{1}{\psi(x) = [\Sigma x_i \varphi'(x'_i)]^2 - 4\varphi(x)\varphi(x').}
If we put $x_i + \lambda x'_i$ in place of $x_i$ in it, then, substituting
for $\varphi(x + \lambda x')$ the expression \S.\,56, (11)
and setting $\Sigma x'_i \varphi'(x'_i) = 2\varphi(x')$, we obtain
\[
\psi(x + \lambda x') = \psi(x),
\]
from which it follows that $\psi(x)$ depends only on $m - 1$ linear combinations
of the $m$ variables. If one puts for $x'_1, x'_2 \ldots x'_m$
arbitrary special values for which $\varphi(x')$ does not vanish, and
takes $x'_1$ different from zero, then
% [sic] The print has x_n; the variables otherwise run through x_m.
\[
x_1'^2 \psi(x_1, x_2, \ldots x_n) = \psi(0, x_2 x'_1 - x_1 x'_2, \ldots x_m x'_1 - x_1 x'_m),
\]
hence a function of the $m - 1$ variables
\[
y_2 = x_2 x'_1 - x_1 x'_2, \ldots y_m = x_m x'_1 - x_1 x'_m.
\]
One can also confirm this by observing that
\[
\psi'(x_1), \ \psi'(x_2) \ldots \psi'(x_m)
\]
all vanish simultaneously for $x_1, x_2 \ldots x_m = x'_1, x'_2 \ldots x'_m$,
so that the determinant of the function $\psi$ must vanish.

Formula (1) now contains a very important result, which
can best be surveyed if we give this formula the following
form.

We set
\weq{2}{\begin{aligned}
\psi(x) &= -4\varphi(x')\chi(y_2, \ldots y_m) \\
\Sigma x_i \varphi'(x'_i) &= 2\sqrt{\varphi(x')}\, Y_1,
\end{aligned}}
and obtain
\weq{3}{\varphi(x) = Y_1^2 + \chi(y_2, y_3 \ldots y_m),}
whereby $\varphi(x)$ is represented as the sum of a square of a
% [sic] The print has lowercase y_1 here, while (2) and (3) use Y_1.
linear function $y_1$ and a function $\chi$ of $m - 1$ variables.

But from this follows the important theorem

\gesperrt{A quadratic form of $m$ variables can always be
represented as a sum of squares of $m$ or fewer
linear functions.}

This theorem is evidently true for a function of one
variable, which is itself a square. But from (3) the truth
of the theorem follows for a function of $m$ variables,
provided its truth is assumed for a function of $m - 1$
variables.
